Data analysts, journalists \cite{young2018makes}, and social scientists \cite{brugnone2024ought, benkler2024recognizing, friedman2024bottom} increasingly seek to turn vast amounts of unstructured text into actionable insights by curating its semantic content. Examples of such corpora include social media discussions, scientific articles, and government white papers. While text summaries can highlight key information like sentiment and argumentation by tying together salient features with narrative flow, modern NLP tools like BERTopic \cite{grootendorst2022bertopic} allow for deeper information extraction and summarization through unsupervised pipelines designed for classification and visualization. The expressiveness of these NLP tools is increasing due to high-quality vector embeddings produced by large language models (LLMs).

Text embeddings are high-dimensional, numerical representations of language, often produced as a byproduct of neural network training. These $d$-dimensional unit vectors live on the $(d-1)$-sphere, where inner products (i.e., cosine similarity) quantify semantic similarity. Recent work \cite{park2024the} suggests that the geometry of text embeddings encodes rich semantic structure, extending beyond cosine similarity.

Dimensionality reduction techniques like Principal Component Analysis (PCA) and Uniform Manifold Approximation and Projection (UMAP) \cite{mcinnes2018umap} aim to capture semantic nuance by leveraging the (linear or nonlinear) manifold structure of embeddings. The manifold hypothesis—asserting that although data lies in a high-dimensional ambient space, it approximately resides on a low-dimensional submanifold \cite{narayanan2010sample}—is central to these approaches. This hypothesis helps circumvent the curse of dimensionality and enables statistical techniques such as kernel density estimation in high dimensions. As demonstrated in, e.g., \citet{coifman2006diffusion}, \citet{moon2019visualizing}, \citet{haghverdi2015diffusion}, and others it is possible to invoke the manifold hypothesis to leverage data geometry and improve upon methods like PCA and t-SNE. Can this perspective be applied to unfurl the expressive geometry hidden within LLM text embeddings? 

\subsection{PHATE}
The manifold hypothesis has inspired significant advances in the analysis of single-cell RNA sequencing (scRNA-seq) data \cite{moon2019visualizing}, which consists of tens-of-thousands of gene expression measurements from thousands of cells. These data exhibit characteristics that include gene dropout and batch effects, which complicate statistical analysis. Moreover, subsamples of cells often follow trajectories along developmental pathways in gene expression space \cite{moon2019visualizing}. These properties led to the development of PHATE, a manifold learning technique that leverages both local and global data geometry to produce low-dimensional embeddings that preserve this structure.

Large-scale textual datasets share many structural characteristics with scRNA-seq data. The uneven expression of social and cultural signals means that important ideas may appear only in small subsets of documents—analogous to dropout. Batch effects also arise in textual data due to differing media preferences (e.g., Bluesky vs. Truth Social) and context dependency, complicating comparison across sources. In addition, narrative arcs in text resemble biological trajectories and may be recoverable through embedding methods. These parallels suggest that PHATE is a promising tool for text analysis. However, this potential has not been systematically evaluated. Here, we present a qualitative analysis of PHATE as a visualization tool and a quantitative assessment of its performance in hierarchical clustering of text.

\subsection{Hierarchical clustering}
A multiscale version of PHATE utilizing diffusion condensation \cite{brugnone2019coarse} has shown promise for identifying hierarchical structure in biological data \cite{kuchroo2022multiscale}. Similar to biology, textual data too contain latent hierarchies that are often overlooked when analyzed at a single level of granularity. Semantic content can be more richly understood by recognizing these hierarchical relationships.

Manually curated resources like WordNet capture lexical hierarchies using synsets. However, unstructured textual corpora typically lack metadata to indicate such relationships. Connections between broad concepts (e.g., “the economy”) and finer-grained topics (e.g., “employment,” “inflation”) must be inferred from the data, which necessitates either laborious and subjective manual labeling or the application of unsupervised machine learning methods.

Early automated approaches inferred hierarchies by analyzing term co-occurrence statistics \cite{sanderson1999deriving}. Modern embeddings have surpassed these methods, enabling more nuanced clustering. The standard pipeline for hierarchical clustering of text involves: (1) embedding the documents, (2) reducing the dimensionality, and (3) applying clustering algorithms at multiple levels of granularity. PCA and UMAP are frequently used for dimensionality reduction, followed by agglomerative clustering or HDBSCAN. Such methods have enabled topic modeling and sentiment analysis, but there may be an opportunity to improve upon them with more expressive coordinates.

In this paper we evaluate PHATE’s ability to generalize to text by examining how well it: (1) recovers known concept hierarchies when coupled with hierarchical clustering, and (2) visualizes semantic relationships among documents. We compare PHATE’s performance to that of PCA, UMAP, and t-SNE (t-distributed Stochastic Neighbor Embedding) across a range of parameter settings and evaluate their ability to produce more informative visualizations and to recover ground-truth hierarchies when coupled with hierarchical clustering.
